\section{The strongest and newest relevant literature}
\label{sec:literature}
\subsection{Direct water-wave evidence}
\paragraph{NLS assimilation and inverse calibration: Ehlers et al. (2024).}
This \emph{Fluids} paper studies 630 synthetic higher-order-spectral (HOS)
wave samples, using observations at two gauge positions and evaluating at
intermediate positions. Fitting carrier-related parameters reduces reported
mean surface similarity parameter (SSP) error from $0.204$ to $0.186$.
Broader spectra and steeper waves remain difficult. The experiments use
numerical data, not a laboratory validation. The paper explicitly proposes
Dysthe residuals, gauge-separation studies, and actual measurements
\cite{ehlers2024}. Its importance is the realistic inverse task and explicit
model discrepancy. Better reconstruction with fitted parameters does not,
by itself, prove that those parameters recover the true physics.

\paragraph{Potential-flow PINNs: Ehlers et al. (2025).}
Published in \emph{Physical Review Fluids} on 2 September 2025, this study
uses elevation and velocity-potential networks with potential-flow
constraints. It includes laboratory flume measurements and reports a
laboratory reconstruction SSP of $0.049$. Forecast quality deteriorates
beyond the assimilation interval; inferred laboratory velocity potential
does not have independent measured ground truth. This is substantial
evidence beyond analytic-soliton fitting, but uses different equations from
the approved envelope model. Pretraining, faster training, and operators
appear in the outlook \cite{ehlers2025}.

\paragraph{Water-wave PINO: Ehlers, Stender and Hoffmann (2025).}
Published in \emph{Physics of Fluids} on 7 October 2025, this paper already
learns sparse-observation reconstruction with an FNO-based physics-informed
operator. It uses 634/211/211 training/validation/test HOS samples and
potential-flow constraints; reported test SSP means are $0.1035$ and
$0.1341$ for two observation configurations. A representative $0.014$\,s
inference time excludes training. These are synthetic, long-crested cases,
not proof of operational ocean forecasting \cite{wavepino}.
Author code is public. A 1,056-sample HOS dataset was released on
8 April 2026 \cite{hosdata}. Availability was checked; neither was run here.
This is the main prior-art obstacle to claiming generic sparse water-wave
operator learning as new.

\subsection{Direct NLS methods to reproduce or compare}
\paragraph{Compact networks and optimization: Urb\'an, Stefanou and Pons (2025).}
The \emph{Journal of Computational Physics} paper is a strong direct
cubic-NLS baseline. Its self-scaled Broyden example uses two hidden layers
of 40 units and reports real/imaginary relative errors of approximately
$(5\pm1,8\pm2)\times10^{-6}$. Hard initial conditions and periodic features
are part of the setup. Its domain differs from the original Raissi
benchmark, expanding from $[-5,5]$ to $[-15,15]$; the nonperiodic
$\operatorname{sech}$ tails become smaller. Consequently the cross-paper
gain cannot be credited to optimization alone. The authors explicitly
suggest further work on modified limited-memory BFGS
\cite{urban}. Reproduce the complete setup before selecting a larger network.

\paragraph{Complex-valued learning: Zhang et al. (2025).}
The \emph{Physical Review Research} CV-PINN paper has particularly useful
accessible evidence: a full manuscript and public implementation. Complex
layers and predictive dynamic monitoring are compared with an approximately
parameter-matched real network across 16 one-dimensional
Gross--Pitaevskii/NLS examples, with three repeats. Some training runs take
approximately two GPU hours. Spectral computations supply reference fields;
this does not show superiority to those solvers. Sampling improvement is
not uniform in every case. The conclusion explicitly proposes richer
residual/solution-gradient sampling \cite{cvpinn}. Treat this as an
architecture and training source, not as water-wave validation.

\paragraph{Hard-constraint ablations: Hao et al. (2026).}
The \emph{Foundations of Data Science} article, early online 8 June 2026,
examines the effect of initial conditions and training stability. Its NLS
combination of hard initial/periodic conditions, self-adaptive weighting,
and mini-batching reports real/imaginary relative errors
$1.80\times10^{-4}$ and $2.58\times10^{-4}$. Code is public
\cite{hardic}. The value lies in ablation and analysis, not the invention of
hard trial functions. The findings do not imply that every noisy inverse
problem should enforce measured initial data exactly. Its numbers should
not be ranked against Urb\'an's numbers without matching the experiments.

\paragraph{Time decomposition and interacting structures: Chen, Li and Yuen (2025).}
Two direct NLS papers study interface-trained time decomposition
(TDD-PINNs) and nonlinear-interaction multi-domain networks (NI-MPINNs).
They include breathers, rogue-wave combinations, and inverse coefficients
with synthetic noise \cite{tdd,interaction}. Their publisher previews report
substantial gains over conventional PINNs. Full comparative budgets and
code were not verified, so the gains are not independently auditable here.
These are relevant precedents for domain decomposition, which should not
be presented as a new idea in a Dysthe application. A spatial formulation
needs propagation-distance slabs, with consistent interface fields.

\paragraph{Transfer between NLS cases: Chen et al. (2024).}
Multi-view transfer learning reuses pretrained NLS information between
related initial/boundary conditions \cite{transfer}. This is a useful
repeated-case comparator, with publisher-level evidence inspected.
Fine-tuning a pretrained solver and learning an input-conditioned operator
are different algorithms. A new study must measure both adaptation cost
and performance on complete unseen wave packets.

\subsection{Very recent 2026 candidates}
\begin{longtable}{@{}P{0.22\textwidth}P{0.35\textwidth}P{0.35\textwidth}@{}}
\toprule
Paper/date & Method and reported evidence & Assessment and access limit \\
\midrule
\endfirsthead
\toprule
Paper/date & Method and reported evidence & Assessment and access limit \\
\midrule
\endhead
GRAVIS, 14 June \cite{gravis} &
Residual magnitude plus output-layer loss-gradient information; seven
localized NLS cases, matched architectures/budgets; reports
$1.56$--$2.37$ times lower errors than RAD. &
Strong recent sampling comparator. Publisher abstract/metadata inspected;
full text, code, and complete cost audit unavailable. Parameter gradients
are not the same as spatial solution gradients. \\
\addlinespace
Noether-constrained PINNs, 5 September issue \cite{noether} &
Dynamic conservation penalties plus adaptive residual sampling; reports
small soliton amplitude error and invariant fluctuations. &
Promising conservation comparator. Abstract-level access; phase,
forecast horizon, repetitions, and matched costs unverified. Soft penalties
do not enforce exact conservation. \\
\addlinespace
CPA-PINNs, 15 October issue \cite{cpa} &
Phase--amplitude representation with continuity, mass and energy penalties;
semiclassical two-dimensional examples and inverse tests with noise. &
Newest verified phase--amplitude candidate here. Publisher preview
inspected; full tables/code unavailable. Defocusing examples do not validate
focusing water-wave instability. \\
\bottomrule
\end{longtable}

These papers justify updating an experimental comparison. They do not
justify calling one method the established best. A separate late watchlist
item is Hu et al.'s \emph{Geometry-Agnostic Data-Driven Discovery of Extreme
Scenarios}, submitted 20 October 2026 \cite{extreme2026}. Its stationary
NLS/coherent-structure examples and proposed extreme-event applications
are not a demonstrated Dysthe forecasting method.

Two WKB-PINN titles were screened but excluded from the confirmed-cutoff
shortlist: DOI \texttt{10.1016/j.cnsns.2026.110653} has a November 2026
issue date without a verified pre-cutoff online date; the exact eligible
date of DOI \texttt{10.1016/j.physd.2026.135394} was not established.
The source ledger retains these uncertainties rather than assigning dates.

\subsection{Transferable methods: useful, with explicit boundaries}
\paragraph{Causality and depth.}
Wang, Sankaran and Perdikaris' causal PINNs weight later evolution stages
according to earlier residuals \cite{causal}. This is not identical to
training separate time windows. PirateNets initializes adaptive residual
gates so that an effectively shallow model gains depth, and combines this
with Fourier embeddings and physics-informed initialization
\cite{pirate}. Neither flagship paper supplied a direct NLS benchmark in
the inspected material. Their application to \eqref{eq:dysthe} is a
testable transfer proposal. PirateNets explicitly points toward
PDE-specific embeddings and operator-learning initialization.

\paragraph{Optimization beyond Adam and L-BFGS.}
The NeurIPS 2025 gradient-alignment paper studies SOAP preconditioning on
ten PDE benchmarks, reporting roughly $2$--$10$ times better accuracy
alongside approximately twice the training expense \cite{soap}.
The inspected text does not establish NLS superiority. Kiyani et al.
provide broader comparisons of self-scaled quasi-Newton methods
\cite{optimizer}. Include a compact quasi-Newton baseline before declaring
that a sophisticated architecture is necessary. Compare wall time to a
target accuracy, not equal optimizer iteration counts.

\paragraph{Operator learning and benchmark discipline.}
Li et al.'s PINO combines data and PDE constraints to learn across input
functions \cite{pino}. That supports the many-query formulation; it is not
evidence for the present Dysthe implementation. PINNacle supplies broader
evidence that method strengths depend on the PDE and problem regime
\cite{pinnacle}. A shared benchmark is more informative than comparing the
best numbers from unrelated papers.

\paragraph{Spectral and goal-oriented residuals are already prior art.}
Spectral informed neural networks replace differential evaluations with
spectral operations \cite{spectral}. Govoeyi and Richter's April 2026
preprint uses dual-weighted-residual ideas for functional-oriented sampling
on Laplace problems \cite{goal}. An operator-budget study explicitly
charges adaptive-sampler residual evaluations \cite{operatorbudget}.
These are relevant precedents; none establishes a solution to the specific
Dysthe nonlocal error-control problem proposed below. Their existence
rules out claiming generic spectral PINNs or generic goal-oriented
sampling as new.

\subsection{Theory and classical numerical controls}
\paragraph{NLS error bounds.}
Alejo, Cossetti, Fanelli, Mu\~noz and Valenzuela study one-dimensional
subcritical focusing NLS on the real line (preprint, revised June 2025)
\cite{alejo}. Their result uses energy/Strichartz-type control, small
defects, and integration assumptions. It is conditional error analysis,
not a guarantee that gradient descent finds a suitable network.
Nonzero-background breathers also need a functional setting different
from directly assuming that the whole field lies in $H^1(\RR)$.
This is the most directly relevant theoretical starting point found here;
it does not immediately cover derivative/nonlocal Dysthe terms.

\paragraph{Residual certification.}
Eiras et al.'s ICML 2024 $\partial$-CROWN bounds worst-case residuals
over continuous domains and includes a Schr\"odinger example
\cite{eiras}. Full text and a public code repository were inspected.
A bound on a residual is not yet a bound on solution or forecast error:
one still needs a stability argument and controlled initial/boundary
defects. The nonlocal operator in \eqref{eq:dysthe} requires additional
analysis beyond a pointwise local residual.

\paragraph{Do not assume PINNs beat numerical solvers.}
Grossmann et al.'s systematic 2024 comparison includes semilinear
Schr\"odinger problems; the tested vanilla PINNs do not outperform FEM in
total solution time and accuracy, though evaluation can sometimes be
faster \cite{grossmann}. This does not rule out newer methods. It does
establish that a neural representation alone is not a competitive
numerical result. For smooth periodic envelopes, include a well-resolved
Fourier-based solver and an efficient tuned solver at each required
accuracy; an excessively expensive reference is not a fair speed baseline.

\clearpage
\subsection{What is actually stated as future work?}
\begin{longtable}{@{}P{0.22\textwidth}P{0.32\textwidth}P{0.38\textwidth}@{}}
\toprule
Source & Explicit direction in inspected text & Consequence for a new paper \\
\midrule
\endfirsthead
\toprule
Source & Explicit direction in inspected text & Consequence for a new paper \\
\midrule
\endhead
Ehlers 2024 \cite{ehlers2024} & Dysthe residual; gauge separation;
actual observations. & A directly motivated extension, requiring a
controlled contribution beyond model substitution. \\
Ehlers 2025 \cite{ehlers2025} & Faster training, pretraining, PINOs. &
The generic operator direction is already pursued by \cite{wavepino}. \\
CV-PINN \cite{cvpinn} & Residual and solution-gradient guided resampling. &
GRAVIS overlaps the broader adaptive-sampling idea; target a specific
physical observable or identifiability mechanism. \\
Urb\'an et al. \cite{urban} & Modified limited-memory BFGS. &
Compare memory and total solve time as well as final loss. \\
PirateNets \cite{pirate} & PDE-specific embeddings; operator initialization. &
Derive features from the approved dispersion relation and test transfer. \\
SOAP \cite{soap} & Less expensive preconditioned algorithms. &
Require matched-accuracy cost improvements, including preconditioner work. \\
\bottomrule
\end{longtable}
The five questions that follow are this review's proposals. They are not
attributed to authors unless explicitly connected above. Directional-wave
and alternative-formulation suggestions in the literature are not part of
the current project plan.
